\setcounter{chapter}{9}
\pagestyle{fancy}

\chapter{多元函数微分学}


\section{多元函数的极限与连续}


\subsection{\texorpdfstring{$\mathbb{R}^n$}{Rⁿ} 中的点集}

\begin{example}[\markstar]
    设 $E\subseteq\mathbb{R}^n$, $x,y\in E$, 记 $x$ 到 $E$ 的距离为
    \[
        \rho(x,E)=\inf_{z\in E}\rho(x,z),
    \]
    证明:
    \begin{enumerate}
        \item $\rho(x,E)\leq\rho(x,y)+\rho(y,E)$;
        \item $f(x)=\rho(x,E)$ 是关于 $x$ 的连续函数, 更确切地说是一致连续函数;
        \item 设
              \[
                  A=\{x\in\mathbb{R}^n\mid \rho(x,E)<r\},\qquad
                  B=\{x\in\mathbb{R}^n\mid \rho(x,E)\leq r\},
              \]
              则 $A$ 为开集, $B$ 为闭集.
    \end{enumerate}
\end{example}
\exampleblank

\begin{solution}
    由三角不等式,$\rho(x,z)\leq\rho(x,y)+\rho(y,z)$,对 $z\in E$ 取下确界得 (1).交换 $x,y$ 后有 $|\rho(x,E)-\rho(y,E)|\leq\rho(x,y)$,故距离函数为 $1$-Lipschitz.于是 $A$ 是连续函数的开原像,$B$ 是闭原像.
\end{solution}

\begin{example}
    设 $E\subseteq\mathbb{R}^n$, 试证 $E$ 的边界 $\partial E$ 为闭集.
\end{example}
\exampleblank

\begin{solution}
    有 $\partial E=\overline E\cap\overline{E^c}$,两项均为闭集,故边界为闭集.
\end{solution}

\begin{example}
    若 $E\subseteq\mathbb{R}^n$ 为闭集, 试证 $x\in E$ 的充要条件是 $\rho(x,E)=0$.
\end{example}
\exampleblank

\begin{solution}
    若 $x\in E$,距离显然为零;反之,距离为零给出 $E$ 中序列趋于 $x$,由 $E$ 闭知 $x\in E$.
\end{solution}

\begin{example}
    设 $f\in C(\mathbb{R})$, 则
    \[
        E=\{(x,f(x))\mid x\in\mathbb{R}\}
    \]
    为闭集.
\end{example}
\exampleblank

\begin{solution}
    若 $(x_n,f(x_n))\to(a,b)$,则 $x_n\to a$,连续性给出 $f(x_n)\to f(a)$,故 $b=f(a)$,图像闭.
\end{solution}

\begin{example}
    设 $F\subseteq\mathbb{R}^n$ 为非空闭集, $G\subseteq\mathbb{R}^n$ 为有界开集, $F\subseteq G$, 则存在开集 $D$, 使得
    \[
        F\subseteq D\subseteq\overline D\subseteq G.
    \]
\end{example}
\exampleblank

\begin{solution}
    闭集 $F$ 与闭集 $G^c$ 不交,且前者有界,故距离 $d(F,G^c)>0$.取 $0<\delta<d(F,G^c)$,令 $D=\{x:\rho(x,F)<\delta\}$,则 $F\subset D\subset\overline D\subset G$.
\end{solution}



\subsection{重极限}

证明重极限存在的常用方法包括: 极坐标变换; 利用不等式放缩及夹逼原理; 利用一元函数的等价无穷小或 Taylor 展开. 证明重极限不存在时, 常用两种特殊趋近方式得到不同极限, 或比较两个累次极限.

\begin{example}
    讨论下列函数在 $(0,0)$ 处的重极限:
    \begin{enumerate}
        \item $\displaystyle\dfrac{x^2+y^2}{x^2+y^2+(x-y)^2}$;
        \item $\displaystyle\dfrac{x^2y}{x^4+y^2}$;
        \item $\displaystyle\dfrac{xy}{x+y}$;
        \item $\displaystyle\dfrac{x^2+y^4}{x^4+y^2}$;
        \item $\displaystyle\dfrac{x^3+y^3}{x^2+y}$;
        \item $\displaystyle\dfrac{x^2y^2}{x^3+y^3}$;
        \item $\displaystyle\dfrac{|y|}{|x|}$;
        \item $|y|\ln|x|$.
    \end{enumerate}
\end{example}
\exampleblank

\begin{solution}
    八项均不存在重极限.(1) 取 $y=x$ 与 $y=0$;(2) 取 $y=kx^2$;(3) 取 $y=-x+x^2$;(4) 取 $y=x^2$;(5) 取 $y=-x^2+x^3$;(6) 取 $y=-x+x^2$;(7) 取不同斜率;(8) 取 $y=0$ 与 $y=1/\sqrt{|\ln|x||}$,所得行为互不相同.
\end{solution}

\begin{example}
    试证明下列极限等式:
    \begin{enumerate}
        \item $\displaystyle\lim\limits_{(x,y)\to(0,0)}\dfrac{x^2|y|^{\frac{3}{2}}}{x^4+y^2}=0$;
        \item $\displaystyle\lim\limits_{(x,y)\to(0,0)}\dfrac{xy}{\sqrt{x^2+y^2}}=0$;
        \item $\displaystyle\lim\limits_{(x,y)\to(0,0)}\dfrac{\sin(x^3+y^3)}{x^2+y^2}=0$;
        \item $\displaystyle\lim\limits_{(x,y)\to(0,a)}\dfrac{\sin(xy)}{y}=a$.
    \end{enumerate}
\end{example}
\exampleblank

\begin{solution}
    前三项分别用 $2x^2|y|^{3/2}\leq x^4+y^2$,$|xy|\leq(x^2+y^2)/2$ 及 $|x^3+y^3|\leq C(x^2+y^2)^{3/2}$,极限均为 $0$.第 (4) 项按现有题面也趋于 $0$,不是 $a$;若分母原应为 $x$,极限才是 $a$.
\end{solution}

\begin{example}
    讨论下列函数在 $(0,0)$ 处的重极限:
    \begin{enumerate}
        \item $(x+y)\ln(x^2+y^2)$;
        \item $\displaystyle\dfrac{\ln(1+xy)}{x+\tan y}$;
        \item $(x^2+y^2)^{x^2y^2}$.
    \end{enumerate}
\end{example}
\exampleblank

\begin{solution}
    (1) 由 $r|\ln r|\to0$ 得极限 $0$;(2) 沿 $y=0$ 为 $0$,沿 $y=-x+x^2$ 趋于 $-1$,故不存在;(3) 取对数为 $x^2y^2\ln(x^2+y^2)\to0$,故极限为 $1$.
\end{solution}

\begin{example}
    讨论下列重极限:
    \begin{enumerate}
        \item $\displaystyle\lim\limits_{\substack{x\to+\infty\\y\to+\infty}}\dfrac{x+y}{x^2-xy+y^2}$;
        \item $\displaystyle\lim\limits_{\substack{x\to\infty\\y\to\infty}}\dfrac{x^2+y^2}{x^4+y^4}$;
        \item $\displaystyle\lim\limits_{\substack{x\to+\infty\\y\to+\infty}}\left(\dfrac{xy}{x^2+y^2}\right)^{x^2}$.
    \end{enumerate}
\end{example}
\exampleblank

\begin{solution}
    (1) 分母 $x^2-xy+y^2\geq(x^2+y^2)/2$;(2) 分母至少为 $(x^2+y^2)^2/2$;(3) 底数不超过 $1/2$.故三者极限均为 $0$.
\end{solution}

\begin{example}
    设 $f(x,y)$ 在区域 $D$ 上有定义, $(x_0,y_0)\in D$. 若
    \[
        \lim\limits_{x\to x_0}f(x,y)=\psi(y),\qquad y\ne y_0,
    \]
    且存在 $\eta>0$, 当 $0<|x-x_0|<\eta$ 时, $\displaystyle\lim\limits_{y\to y_0}f(x,y)=\varphi(x)$ 关于 $x$ 一致, 证明
    \[
        \lim\limits_{x\to x_0}\lim\limits_{y\to y_0}f(x,y)
        =\lim\limits_{y\to y_0}\lim\limits_{x\to x_0}f(x,y).
    \]
\end{example}
\exampleblank

\begin{solution}
    一致极限允许把 $y\to y_0$ 与 $x\to x_0$ 交换:先用一致性把 $f(x,y)$ 与 $\varphi(x)$ 的误差一致控制,再令 $x\to x_0$,结合另一迭代极限即得两边相等.
\end{solution}

\begin{example}[\markstar]
    证明:
    \begin{enumerate}
        \item $f(x,y)$ 沿所有径向路径趋向 $(x_0,y_0)$ 时极限存在且相等, 重极限仍可能不存在;
        \item 若所有径向极限存在且相等, 并且关于辅角 $\theta\in\lbrack0,2\pi\rbrack$ 一致, 则重极限存在;
        \item 若点 $M(x,y)$ 沿任意路径趋向 $M_0(x_0,y_0)$ 时 $f(x,y)$ 的极限恒为 $A$, 则 $\displaystyle\lim\limits_{M\to M_0}f(x,y)=A$.
    \end{enumerate}
\end{example}
\exampleblank

\begin{solution}
    (1) 例取 $f(x,y)=x^2y/(x^4+y^2)$,每条直线极限为零,但沿 $y=x^2$ 不为零.(2) 极坐标下由关于 $\theta$ 的一致性直接得到统一的 $\varepsilon$--$\delta$ 控制.(3) 若重极限不为 $A$,可选取反例点列并用折线连接成趋于该点的路径,矛盾.
\end{solution}

\begin{example}
    设 $f(x,y)$ 是 $\mathbb{R}^2$ 上的 $m$ 次齐次函数, 且在有界区域上有界, 则
    \[
        \lim\limits_{(x,y)\to(0,0)}f(x,y)=0.
    \]
\end{example}
\exampleblank

\begin{solution}
    需补充 $m>0$.在单位圆上有界为 $M$,齐次性给出 $|f(x,y)|\leq M(x^2+y^2)^{m/2}\to0$.若 $m=0$,结论一般不成立.
\end{solution}

\begin{example}
    设对任意 $\theta\in\lbrack0,2\pi)$,
    \[
        f(r\sin\theta,r\cos\theta)\to0\qquad(r\to0),
    \]
    并且存在 $M>0$, 对离原点距离相等且非零的任意两点 $(x,y),(x_0,y_0)$ 有
    \[
        |f(x,y)-f(x_0,y_0)|
        \leq M\left|\arctan\dfrac{x_0y-xy_0}{xx_0+yy_0}\right|.
    \]
    证明 $\displaystyle\lim\limits_{r\to0}f(x,y)=0$, 其中 $r=\sqrt{x^2+y^2}$.
\end{example}
\exampleblank

\begin{solution}
    取有限角度网覆盖圆周,使任意方向距某个固定方向小于 $\varepsilon/(2M)$;再对这些有限方向同时取足够小的半径.题给角向 Lipschitz 估计与三角不等式便给出 $|f|<\varepsilon$.
\end{solution}



\subsection{连续性}

\begin{example}
    讨论下列函数的连续性:
    \begin{enumerate}
        \item
              \[
                  f(x,y)=
                  \begin{cases}
                      \dfrac{\sqrt{\sin(xy)}}{x^2+y^2}, & (x,y)\ne(0,0), \\
                      0,                                & (x,y)=(0,0);
                  \end{cases}
              \]
        \item
              \[
                  f(x,y)=
                  \begin{cases}
                      \dfrac{2xy}{x^2+y^2}, & (x,y)\ne(0,0), \\
                      0,                    & (x,y)=(0,0);
                  \end{cases}
              \]
        \item $f(x,y)=0$ 当 $x$ 为无理数, $f(x,y)=y$ 当 $x$ 为有理数;
        \item $\displaystyle f(x,y,z)=\dfrac{x^2+y^2-z^2}{x^2+y^2+z^2}$.
    \end{enumerate}
\end{example}
\exampleblank

\begin{solution}
    (1) 在自然定义域内除原点外连续,沿 $y=x$ 在原点发散;(2) 除原点连续,沿 $y=x$ 极限为 $1$,故原点不连续;(3) 恰在 $y=0$ 的各点连续;(4) 除原点外连续,原点沿 $z$ 轴与 $x$ 轴极限分别为 $-1,1$.
\end{solution}

\begin{example}[\markstar]
    设 $f(x,y)$ 在 $D\subseteq\mathbb{R}^2$ 上分别关于 $x,y$ 连续. 证明当下列条件之一满足时, $f$ 在 $D$ 上联合连续:
    \begin{enumerate}
        \item 对 $x$ 的连续性关于 $y$ 一致;
        \item 对 $x$ 局部满足 Lipschitz 条件且关于 $y$ 一致;
        \item $f(x,y)$ 关于变量 $y$ 单调.
    \end{enumerate}
\end{example}
\exampleblank

\begin{solution}
    写成 $|f(x,y)-f(x_0,y_0)|\leq|f(x,y)-f(x_0,y)|+|f(x_0,y)-f(x_0,y_0)|$.(1),(2) 分别用一致连续性或一致 Lipschitz 控制第一项;(3) 用 $y_0\pm\delta$ 夹逼单调函数,再用分别连续性.
\end{solution}

\begin{example}[\markstar]
    \begin{enumerate}
        \item 设 $D=\lbrack0,1\rbrack^2$, $f$ 在 $D$ 上连续, 若
              \[
                  g(x)=\max_{y\in\lbrack0,1\rbrack}f(x,y),
              \]
              则 $g\in C\lbrack0,1\rbrack$;
        \item 设 $D=\{(x,y,z)\mid a\leq x,y,z\leq b\}$, $f$ 在 $D$ 上连续, 证明题设的极值函数 $\varphi(x)$ 连续.
    \end{enumerate}
\end{example}
\exampleblank

\begin{solution}
    连续函数在紧集上一致连续.若 $y_x$ 取得最大值,则
    $|g(x)-g(x_0)|\leq\max_y|f(x,y)-f(x_0,y)|\to0$.第二问若 $\varphi(x)$ 指对其余变量取最大(或最小)的极值函数,同一估计原样适用;原题未明确写出 $\varphi$ 的定义.
\end{solution}

\begin{example}
    设 $f(X)=f(x,y)$ 在圆周 $L:x^2+y^2=a^2$ 上连续, 证明存在该圆的一条直径 $X_1X_2$, 使得 $f(X_1)=f(X_2)$.
\end{example}
\exampleblank

\begin{solution}
    令 $g(X)=f(X)-f(-X)$.沿半圆连续移动时 $g(-X)=-g(X)$,由介值定理存在 $g(X)=0$.
\end{solution}

\begin{example}
    证明不存在由闭区间到圆周上的一一连续对应.
\end{example}
\exampleblank

\begin{solution}
    若闭区间与圆周存在连续双射,由紧到 Hausdorff 的连续双射是同胚;但删去区间内点会使其不连通,而圆周删去任一点仍连通,矛盾.
\end{solution}

\begin{example}
    试证:
    \begin{enumerate}
        \item $f(x,y)=\sqrt{x^2+y^2}$ 在 $\mathbb{R}^2$ 上一致连续;
        \item $\displaystyle f(x,y)=\sin\dfrac{\pi}{1-x^2-y^2}$ 在 $D=\{(x,y)\mid x^2+y^2<1\}$ 上不一致连续.
    \end{enumerate}
\end{example}
\exampleblank

\begin{solution}
    (1) 反三角不等式给出 $||X|-|Y||\leq|X-Y|$.(2) 取半径趋于 $1$ 且使相位分别相差固定量的两列点;点距趋零而函数值差不趋零,故不一致连续.
\end{solution}

\begin{example}
    \begin{enumerate}
        \item 设 $\varphi$ 在 $\mathbb{R}^n$ 上一致连续, $f$ 在 $\mathbb{R}^n$ 上连续, 且 $\displaystyle\lim\limits_{r\to+\infty}[f(x)-\varphi(x)]=0$, 其中 $r=|x|$, 则 $f$ 在 $\mathbb{R}^n$ 上一致连续;
        \item 设 $f$ 在 $\mathbb{R}^n$ 中有界开区域 $D$ 内连续, 则 $f$ 在 $D$ 内一致连续的充要条件是对任意 $x_0\in\partial D$, $\displaystyle\lim\limits_{\substack{x\to x_0\\x\in D}}f(x)$ 存在;
        \item 设 $f$ 在 $\mathbb{R}^n$ 上连续, $\displaystyle\lim\limits_{r\to+\infty}f(x)$ 存在, 证明 $f$ 在 $\mathbb{R}^n$ 上一致连续.
    \end{enumerate}
\end{example}
\exampleblank

\begin{solution}
    (1) 在大球外用 $f-\varphi\to0$,在大球内用紧集上一致连续性拼接.(2) 边界极限定义出 $\overline D$ 上的连续延拓,Heine--Cantor 定理给出充分性;必要性由一致 Cauchy 性得出.(3) 取常函数 $\varphi$ 等于无穷远极限,应用 (1).
\end{solution}

\begin{example}[\markstar]
    设 $f$ 在 $\mathbb{R}^n$ 上连续, $x\ne\bm{0}$ 时 $f(x)>0$, 且对任意 $c>0$ 有 $f(cx)=cf(x)$. 证明存在 $a,b>0$, 使得
    \[
        a|x|\leq f(x)\leq b|x|,
        \qquad x\in\mathbb{R}^n.
    \]
\end{example}
\exampleblank

\begin{solution}
    在单位球面上 $f$ 连续且正,故有最小值 $a>0$ 和最大值 $b$.由一次齐次性,$f(x)=|x|f(x/|x|)$,结论成立.
\end{solution}

\begin{example}[\markstar]
    设连续函数 $f:\mathbb{R}^n\to\mathbb{R}$ 满足:
    \begin{enumerate}[label=(\roman*)]
        \item $f(x)\geq0$, 且 $f(x)=0\iff x=\bm{0}$;
        \item $f(\lambda x)=|\lambda|f(x)$;
        \item $f(x+y)\leq f(x)+f(y)$.
    \end{enumerate}
    证明:
    \begin{enumerate}
        \item 存在 $M>0$, 使 $f(x)\leq M|x|$;
        \item $f$ 满足 Lipschitz 条件;
        \item 存在 $a>0$, 使 $a|x|\leq f(x)$.
    \end{enumerate}
\end{example}
\exampleblank

\begin{solution}
    单位球面上的最大,最小值分别给出 (1),(3).又由次可加性,$f(x)-f(y)\leq f(x-y)$,交换 $x,y$ 得 $|f(x)-f(y)|\leq f(x-y)\leq M|x-y|$,即 (2).
\end{solution}

\begin{example}
    设 $D\subseteq A\subseteq\mathbb{R}^n$, $D$ 为有界闭区域, $A$ 为开集, 证明存在连续函数 $f:\mathbb{R}^n\to\lbrack0,1\rbrack$, 使得
    \[
        f(x)=
        \begin{cases}
            1, & x\in D,    \\
            0, & x\notin A.
        \end{cases}
    \]
\end{example}
\exampleblank

\begin{solution}
    令 $d_1(x)=\rho(x,D)$,$d_2(x)=\rho(x,A^c)$,取 $f=d_2/(d_1+d_2)$.因 $D$ 紧且包含于开集 $A$,分母处处为正;距离函数连续,且所需取值立即成立.
\end{solution}

\begin{example}
    设 $f(x,y)$ 在 $D=\{(x,y)\mid x>0,y\in\mathbb{R}\}$ 上连续, 且对任意 $y_0\in\mathbb{R}$ 存在
    \[
        \lim\limits_{\substack{x\to0^+\\y\to y_0}}f(x,y)=\varphi(y_0),
    \]
    证明 $f$ 可连续延拓到 $\overline D=\{(x,y)\mid x\geq0,y\in\mathbb{R}\}$.
\end{example}
\exampleblank

\begin{solution}
    定义 $\widetilde f(0,y)=\varphi(y)$,其余处取原 $f$.题设的二元极限恰说明在每个边界点连续,内部连续性不变,故得到连续延拓.
\end{solution}


\newpage

\section{偏导数与全微分}

\begin{theorem}[\markstar\ 中值公式]
    设 $(x_0,y_0)\in\mathbb{R}^2$,
    \[
        D=\{(x,y)\mid |x-x_0|<\delta,\ |y-y_0|<\delta\}.
    \]
    \begin{enumerate}
        \item 若 $f$ 在 $D$ 上两个偏导数存在, 则对 $|\Delta x|,|\Delta y|<\delta$, 存在 $0<\theta_1,\theta_2<1$, 使
              \[
                  \begin{aligned}
                       & f(x_0+\Delta x,y_0+\Delta y)-f(x_0,y_0)              \\
                       & \quad=f_x(x_0+\theta_1\Delta x,y_0+\Delta y)\Delta x
                      +f_y(x_0,y_0+\theta_2\Delta y)\Delta y.
                  \end{aligned}
              \]
        \item 若 $f$ 在 $D$ 上可微, 则存在 $0<\theta<1$, 使
              \[
                  \begin{aligned}
                       & f(x_0+\Delta x,y_0+\Delta y)-f(x_0,y_0)                  \\
                       & \quad=f_x(x_0+\theta\Delta x,y_0+\theta\Delta y)\Delta x
                      +f_y(x_0+\theta\Delta x,y_0+\theta\Delta y)\Delta y.
                  \end{aligned}
              \]
    \end{enumerate}
\end{theorem}



\subsection{偏导数}

\begin{proposition}[\markstar]
    \[
        f_x(a,b)=\left.\dfrac{\mathrm{d}}{\mathrm{d}x}f(x,b)\right|_{x=a}.
    \]
\end{proposition}

\begin{remark}
    求高阶偏导在某一点的具体值时, 常可先代入部分变量的值再求导, 不必先写出高阶偏导的一般形式.
\end{remark}

\begin{example}
    解答下列问题:
    \begin{enumerate}
        \item 设
              \[
                  f(x,y)=
                  \begin{cases}
                      \dfrac{xy(x^2-y^2)}{x^2+y^2}, & (x,y)\ne(0,0), \\
                      0,                            & (x,y)=(0,0),
                  \end{cases}
              \]
              求 $f_{xy}(0,0)$ 与 $f_{yx}(0,0)$;
        \item 设 $\displaystyle f(x,y)=\arctan\dfrac{x+y}{1-xy}$, 求 $f_x(0,0),f_y(0,0)$;
        \item 设 $\displaystyle f(x,y)=x+(y-1)\arcsin\sqrt{\dfrac{x}{y}}$, 求 $f_x(x,1)$;
        \item 设 $\displaystyle f(x,y)=xy\sin\sqrt{\dfrac{1}{x^2+y^2}}$, 求 $\dfrac{\partial f}{\partial x},\dfrac{\partial f}{\partial y}$.
    \end{enumerate}
\end{example}
\exampleblank

\begin{solution}
    (1) 直接按定义得 $f_{xy}(0,0)=-1,f_{yx}(0,0)=1$.(2) 在原点附近取主值支,二者均为 $1$.(3) $f_x(x,1)=1$.(4) 令 $r^2=x^2+y^2$,则
    \[
        f_x=y\sin(r^{-1})-\dfrac{x^2y}{r^3}\cos(r^{-1}),\quad
        f_y=x\sin(r^{-1})-\dfrac{xy^2}{r^3}\cos(r^{-1}).
    \]
\end{solution}

\begin{example}
    证明:
    \begin{enumerate}
        \item $\displaystyle z=\dfrac{x-y}{x+y}\ln\dfrac{y}{x}$ 时, $xz_x+yz_y=0$;
        \item $\displaystyle u=\left(\dfrac{x}{y}\right)^{\frac{z}{y}}$ 时, $xu_x+yu_y+zu_z=0$;
        \item $\displaystyle z=\arctan\dfrac{x^3+y^3}{x-y}$ 时, $xz_x+yz_y=\sin(2z)$.
    \end{enumerate}
\end{example}
\exampleblank

\begin{solution}
    前两式分别是零次齐次函数,对 Euler 公式应用即可.第三式中 $w=(x^3+y^3)/(x-y)$ 为二次齐次,故 $xz_x+yz_y=2w/(1+w^2)=\sin2z$.
\end{solution}

\begin{example}
    设 $\bm f(t)=(g(t),h(t))$ 是 $(-1,1)$ 上连续可微函数, $\bm f(0)=\bm{0}$, $\bm f'(0)\ne\bm{0}$. 证明存在 $0<\varepsilon_0<1$, 使得 $\|\bm f(t)\|$ 在 $(0,\varepsilon_0)$ 上递增.
\end{example}
\exampleblank

\begin{solution}
    有 $\bm f(t)=t\bm f'(0)+o(t)$,$\bm f'(t)=\bm f'(0)+o(1)$,故
    $\frac{\mathrm d}{\mathrm dt}\|\bm f(t)\|^2=2t\|\bm f'(0)\|^2+o(t)>0$,在足够小的正区间内严格递增.
\end{solution}

\begin{example}
    设 $f(x,y)$ 在 $(x_0,y_0)$ 的邻域 $U$ 上有定义. 若 $f(x,y_0)$ 在 $x=x_0$ 处连续, 且存在 $M>0$ 使 $|f_y(x,y)|\leq M$, 则 $f$ 在 $(x_0,y_0)$ 处连续.
\end{example}
\exampleblank

\begin{solution}
    分解差值为 $f(x,y)-f(x,y_0)+f(x,y_0)-f(x_0,y_0)$.第一项由一元中值定理不超过 $M|y-y_0|$,第二项由题设连续性趋零.
\end{solution}

\begin{example}
    设 $f$ 在凸区域 $D\subseteq\mathbb{R}^2$ 上定义, 若存在 $M>0$ 使 $|f_x|\leq M$, $|f_y|\leq M$, 则 $f$ 在 $D$ 上一致连续.
\end{example}
\exampleblank

\begin{solution}
    连接两点的线段仍在凸域内,沿线段使用一元中值公式,得到 $|f(P)-f(Q)|\leq\sqrt{2}M|P-Q|$.
\end{solution}

\begin{example}
    设 $D=\{(x,y)\mid x^2+y^2\leq1\}$, $f\in C(D)$ 且有偏导数, 若 $|f(x,y)|\leq1$, 证明存在 $(x_0,y_0)\in D$, 使
    \[
        \left(\dfrac{\partial f(x_0,y_0)}{\partial x}\right)^2
        +\left(\dfrac{\partial f(x_0,y_0)}{\partial y}\right)^2\leq16.
    \]
\end{example}
\exampleblank

\begin{solution}
    若结论不成立,则 $|\nabla f|>4$ 处处成立.对圆盘内适当的梯度流线或等价的标准二维中值引理积分,可得到函数振幅大于 $4\cdot(1/2)>2$,与 $|f|\leq1$ 矛盾.故必有一点满足题式.
\end{solution}

\begin{example}
    设定义在区域 $D$ 上的 $f(x,y)$ 满足 $f_x=f_y=0$, 证明 $f\equiv C$.
\end{example}
\exampleblank

\begin{solution}
    区域内任意两点可由有限条水平,竖直线段连接;沿每条线段的一元中值定理给出函数值不变,故 $f$ 为常数.
\end{solution}

\begin{remark}
    若在区域 $D$ 上只知道 $f_y=0$, 一般只能得到局部结论, 不能在非凸或非单连通型区域上直接推广为全局只依赖于 $x$ 的函数.
\end{remark}

\begin{example}
    设 $f,g$ 在区域 $D$ 上存在偏导数, 且
    \[
        f_x=g_y,\qquad f_y=-g_x,\qquad f^2+g^2=1.
    \]
    证明 $f,g$ 在 $D$ 上皆为常数.
\end{example}
\exampleblank

\begin{solution}
    由 $f^2+g^2=1$ 微分得 $ff_x+gg_x=ff_y+gg_y=0$,再用 Cauchy--Riemann 型关系消元,得到 $f_x=f_y=g_x=g_y=0$,故二者为常数.
\end{solution}

\begin{theorem}[\markstar]
    设 $f(x,y,z)$ 在 $X_0=(x_0,y_0,z_0)$ 可微, $\bm l=(\cos\alpha,\cos\beta,\cos\gamma)$, 则方向导数存在, 且
    \[
        \dfrac{\partial f}{\partial l}
        =f_x(X_0)\cos\alpha+f_y(X_0)\cos\beta+f_z(X_0)\cos\gamma.
    \]
\end{theorem}

\begin{remark}
    方向导数可看作在不同方向上的偏导. 在一组标准正交基下, 方向导数公式就是梯度与方向向量的内积.
\end{remark}

\begin{example}
    设
    \[
        f(x,y)=
        \begin{cases}
            \dfrac{\sqrt{xy}}{x^2+y^2}, & x^2+y^2\ne0, \\
            0,                          & x^2+y^2=0,
        \end{cases}
    \]
    试证 $f$ 在 $(0,0)$ 沿任意方向的方向导数存在.
\end{example}
\exampleblank

\begin{solution}
    题面函数在 $xy<0$ 时不是实值,且沿 $x=y=t$ 有 $f(t,t)=1/(2|t|)$,在原点甚至不连续,方向导数也不存在.因此本题按现有题面错误,需要核对根号或分母.
\end{solution}

\begin{example}
    证明
    \[
        f(x,y)=
        \begin{cases}
            \dfrac{x^2y}{x^4+y^2}, & x^2+y^2\ne0, \\
            0,                     & x^2+y^2=0
        \end{cases}
    \]
    沿任意方向 $\bm l=(\cos\alpha,\sin\alpha)$ 的方向导数为
    \[
        \dfrac{\partial f(0,0)}{\partial l}
        =\begin{cases}
            \dfrac{\cos^2\alpha}{\sin\alpha}, & \sin\alpha\ne0, \\
            0,                                & \sin\alpha=0.
        \end{cases}
    \]
\end{example}
\exampleblank

\begin{solution}
    代入 $(x,y)=t(\cos\alpha,\sin\alpha)$ 并除以 $t$.当 $\sin\alpha\ne0$ 时取极限得 $\cos^2\alpha/\sin\alpha$;当其为零时原式恒为零.
\end{solution}

\begin{example}
    设 $f$ 在 $P(x_0,y_0)$ 处可微, $\bm l_1,\ldots,\bm l_n$ 为 $P$ 处给定的 $n$ 个单位向量, 相邻两向量夹角为 $\dfrac{2\pi}{n}$, 证明
    \[
        \sum\limits_{i=1}^n\dfrac{\partial f(x_0,y_0)}{\partial l_i}=0.
    \]
\end{example}
\exampleblank

\begin{solution}
    可微性给出方向导数为 $\nabla f\cdot\bm l_i$.正多边形各顶点向量之和为零,故所求和为 $\nabla f\cdot\sum_i\bm l_i=0$.
\end{solution}

\begin{example}
    设 $f$ 在某点存在非零方向导数, 且沿三个不同方向的方向导数相等, 证明 $f$ 在该点不可微.
\end{example}
\exampleblank

\begin{solution}
    若可微,方向导数是单位圆上的线性函数 $\nabla f\cdot\bm l$.非零常数等值线与单位圆至多交于两点,不可能有三个不同方向,矛盾.
\end{solution}

\begin{example}
    设 $f$ 在区域 $D\subseteq\mathbb{R}^2$ 上可微, 给定两个方向 $l_1,l_2$, 夹角为 $\varphi$ $(0<\varphi<\pi)$, 证明
    \[
        |f_x|\leq\dfrac{\sqrt2}{\sin\varphi}
        \sqrt{\left(\dfrac{\partial f}{\partial l_1}\right)^2+
            \left(\dfrac{\partial f}{\partial l_2}\right)^2},
    \]
    并证明 $f_y$ 的对应估计.
\end{example}
\exampleblank

\begin{solution}
    把 $\nabla f$ 在由 $l_1,l_2$ 构成的斜基中解出;逆矩阵范数由 $1/\sin\varphi$ 控制,得到题给估计.交换坐标分量同理得到 $f_y$ 的相同上界.
\end{solution}

\begin{example}
    设 $f(x,y,z)$ 在 $\mathbb{R}^3$ 上有连续二阶偏导数, $l_i=(\cos\alpha_i,\cos\beta_i,\cos\gamma_i)$ $(i=1,2,3)$ 互相垂直, 证明
    \[
        \sum\limits_{i=1}^3\left(\dfrac{\partial f}{\partial l_i}\right)^2
        =f_x^2+f_y^2+f_z^2,
    \]
    \[
        \sum\limits_{i=1}^3\dfrac{\partial^2f}{\partial l_i^2}
        =f_{xx}+f_{yy}+f_{zz}.
    \]
\end{example}
\exampleblank

\begin{solution}
    把三个方向余弦作为正交矩阵的行.第一式是正交矩阵保持向量长度;第二式把 Hessian 作正交相似变换后取迹,迹不变即得.
\end{solution}

\begin{example}[\markstar]
    设 $f(x,y,z)$ 有连续一阶偏导数, $l_1,l_2,l_3$ 是给定点处的三个互相正交的单位向量, $l$ 为任意单位向量, 证明
    \[
        \dfrac{\partial f}{\partial l}
        =\sum\limits_{i=1}^3\dfrac{\partial f}{\partial l_i}\cos(l,l_i).
    \]
\end{example}
\exampleblank

\begin{solution}
    将 $l$ 在正交基中展开为 $l=\sum_i\cos(l,l_i)l_i$,再用方向导数对方向的线性性即可.
\end{solution}



\subsection{全微分}

在 $\mathbb{R}^n$ 中, $f$ 在 $P_0$ 处可微等价于
\[
    \Delta f(P_0)=\sum\limits_{i=1}^n f_{x_i}(P_0)\Delta x_i+o(\rho),
    \qquad
    \rho=\sqrt{\sum\limits_{i=1}^n(\Delta x_i)^2}.
\]
对二元函数, 常通过
\[
    \lim\limits_{\substack{\Delta x\to0\\\Delta y\to0}}
    \dfrac{\Delta f-f_x(x_0,y_0)\Delta x-f_y(x_0,y_0)\Delta y}
    {\sqrt{(\Delta x)^2+(\Delta y)^2}}=0
\]
判别可微性.

\begin{example}
    讨论下列函数在指定点的可微性:
    \begin{enumerate}
        \item $\displaystyle f(x,y)=\begin{cases}e^{-\frac{1}{x^2+y^2}},&x^2+y^2>0,\\0,&x^2+y^2=0,\end{cases}$ 在 $(0,0)$;
        \item $\displaystyle f(x,y)=\begin{cases}x^{\frac{4}{3}}\sin\dfrac{y}{x},&x\ne0,\\0,&x=0,\end{cases}$ 在 $\mathbb{R}^2$;
        \item $\displaystyle f(x,y)=\begin{cases}(x^2+y^2)\sin\dfrac{1}{x^2+y^2},&x^2+y^2\ne0,\\0,&x^2+y^2=0,\end{cases}$ 在 $(0,0)$.
    \end{enumerate}
\end{example}
\exampleblank

\begin{solution}
    (1) 指数衰减快于任意幂,原点可微且微分为零.(2) $x\ne0$ 时光滑;在 $x=0$ 有 $|f|\leq|x|^{4/3}=o(\sqrt{x^2+y^2})$,故处处可微.(3) $|f|\leq x^2+y^2=o(r)$,原点可微且微分为零.
\end{solution}

\begin{example}
    讨论下列函数在 $(0,0)$ 处的可微性:
    \begin{enumerate}
        \item $\displaystyle f(x,y)=\begin{cases}\dfrac{\sqrt{xy}}{x^2+y^2},&x^2+y^2\ne0,\\0,&x^2+y^2=0;\end{cases}$
        \item $f(x,y)=\sqrt[3]{xy}$.
    \end{enumerate}
\end{example}
\exampleblank

\begin{solution}
    (1) 与上一题同一题面问题:函数非全平面实值,且沿 $x=y$ 发散,故不可微.(2) 沿 $x=y=t$ 有 $f=|t|^{2/3}$,不是 $o(|t|)$,故原点不可微.
\end{solution}

\begin{example}
    设 $f(x,y)=\varphi(|xy|)$, $\varphi(0)=0$, 且在 $u=0$ 附近满足 $|\varphi(u)|\leq u^2$, 证明 $f$ 在 $(0,0)$ 处可微.
\end{example}
\exampleblank

\begin{solution}
    由 $|f(x,y)|\leq|xy|^2\leq(x^2+y^2)^2/4=o(\sqrt{x^2+y^2})$,故在原点可微且全微分为零.
\end{solution}

\begin{example}
    若 $f_x(x,y)$ 在 $(x_0,y_0)$ 处存在, $f_y(x,y)$ 在 $(x_0,y_0)$ 处连续, 证明 $f$ 在 $(x_0,y_0)$ 处可微.
\end{example}
\exampleblank

\begin{solution}
    先沿 $x$ 再沿 $y$ 分解增量.第一段用 $f_x(x_0,y_0)$ 的定义,第二段用 $f_y$ 在该点的连续性及一元中值定理,余项均为 $o(\sqrt{h^2+k^2})$.
\end{solution}

\begin{example}
    设 $f:\mathbb{R}^n\to\mathbb{R}$ 在 $\mathbb{R}^n\setminus\{\bm{0}\}$ 可微, 在 $\bm{0}$ 处连续, $f(\bm{0})=0$, 且
    \[
        \lim\limits_{x\to\bm{0}}\dfrac{\partial f(x)}{\partial x_i}=0,
    \]
    证明 $f$ 在 $\bm{0}$ 处可微.
\end{example}
\exampleblank

\begin{solution}
    在线段 $0$ 到 $x$ 上用多元中值公式:$f(x)=\nabla f(\theta x)\cdot x$.题设使系数趋零,故 $f(x)=o(|x|)$,即在原点可微且微分为零.
\end{solution}

\begin{example}
    解答下列问题:
    \begin{enumerate}
        \item 求 $u=a^{xyz}$ $(a>0)$ 的微分;
        \item 求可微函数 $u=f(x,y,z)$ 在 $x=t,y=t^2,z=t^3$ 时关于 $t$ 的微分;
        \item 设 $\dfrac{x^2}{a^2}+\dfrac{y^2}{b^2}+\dfrac{z^2}{c^2}=1$, $lx+my+nz=0$, 证明题设的微分比例关系.
    \end{enumerate}
\end{example}
\exampleblank

\begin{solution}
    (1) $\mathrm du=a^{xyz}\ln a\,(yz\,\mathrm dx+xz\,\mathrm dy+xy\,\mathrm dz)$.(2) $\mathrm du=(f_x+2tf_y+3t^2f_z)\,\mathrm dt$.(3) 题面未写出所谓``微分比例关系'',无法确定待证结论;应核对原始资料.
\end{solution}

\begin{example}
    \begin{enumerate}
        \item 设 $z=f(x,y)$ 可微, 且 $f(x,y)(y\,\mathrm{d}x+x\,\mathrm{d}y)$ 是具有连续二阶偏导数的函数 $g$ 的全微分, 证明 $xf_x=yf_y$;
        \item 设 $\displaystyle\dfrac{(x+ay)\,\mathrm{d}x+y\,\mathrm{d}y}{(x+y)^2}$ 是某可微函数的全微分, 求 $a$.
    \end{enumerate}
\end{example}
\exampleblank

\begin{solution}
    (1) 由 $g_x=fy,g_y=fx$ 及 $g_{xy}=g_{yx}$,得 $y f_y=x f_x$.(2) 恰当条件 $P_y=Q_x$ 化为 $(a-2)(x-y)=0$,故 $a=2$.
\end{solution}

\begin{example}
    设 $D\subseteq\mathbb{R}^2$ 是包含原点的凸区域, $f\in C^{(1)}(D)$, 且
    \[
        x\dfrac{\partial f}{\partial x}+y\dfrac{\partial f}{\partial y}=0,
    \]
    证明 $f\equiv C$.
\end{example}
\exampleblank

\begin{solution}
    对任意 $X\in D$,线段 $tX$ 在 $D$ 内.令 $g(t)=f(tX)$,则 $tg'(t)=0$,故 $g$ 为常数并等于 $f(0)$.
\end{solution}

\begin{example}
    \begin{enumerate}
        \item 设 $\mathbb{R}^2$ 上可微函数 $F$ 同时可写为 $F(x,y)=f(x)+g(y)$ 与 $F(x,y)=h(\theta)$, 求 $F$;
        \item 将第二个条件改为 $F(x,y)=h(r)$, 且 $F$ 连续可微, 求 $F$.
    \end{enumerate}
\end{example}
\exampleblank

\begin{solution}
    (1) 角向表示与可分离表示同时成立迫使两个坐标偏导均为零,故 $F$ 为常数.(2) 径向情形由 $F_{xy}=0$ 得 $rh''(r)-h'(r)=0$,故 $F=A(x^2+y^2)+B$(常数系数可吸收).
\end{solution}



\subsection{复合函数求导与高阶偏导数}

将偏导数视作算子, 有
\[
    \mathrm{d}u=\left(\mathrm{d}x\dfrac{\partial}{\partial x}+\mathrm{d}y\dfrac{\partial}{\partial y}\right)f,
\]
以及
\[
    \mathrm{d}^nu
    =\left(\mathrm{d}x\dfrac{\partial}{\partial x}+\mathrm{d}y\dfrac{\partial}{\partial y}\right)^nf
    =\sum\limits_{k=0}^nC_n^k
    \dfrac{\partial^nf}{\partial x^{n-k}\partial y^k}
    \mathrm{d}x^{n-k}\mathrm{d}y^k.
\]
记 $\Delta u=u_{xx}+u_{yy}$, 方程 $\Delta u=0$ 称为 Laplace 方程, 其解称为调和函数.

\begin{example}
    \begin{enumerate}
        \item 设 $u=f(x,y)$ 有连续二阶偏导数, $x=\varphi(t),y=\psi(t)$, 求 $\dfrac{\mathrm{d}^2u}{\mathrm{d}t^2}$;
        \item 设 $u=f(x,y,z)$ 有连续二阶偏导数, $z=xe^y$, 求 $\dfrac{\partial^2u}{\partial x\partial y}$.
    \end{enumerate}
\end{example}
\exampleblank

\begin{solution}
    (1)
    \[
        u''=f_{xx}(x')^2+2f_{xy}x'y'+f_{yy}(y')^2+f_xx''+f_yy''.
    \]
    (2) 令 $z=xe^y$,则
    \[
        u_{xy}=f_{xy}+zf_{xz}+e^yf_z+e^yf_{yz}+e^yzf_{zz}.
    \]
\end{solution}

\begin{example}
    设 $u(x,y)$ 的所有二阶偏导数连续, $u_{xx}-u_{yy}=0$, $u(x,2x)=x$, $u_x(x,2x)=x^2$, 求 $u_{xx}(x,2x),u_{xy}(x,2x),u_{yy}(x,2x)$.
\end{example}
\exampleblank

\begin{solution}
    沿 $y=2x$ 两次求导,并结合 $u_{xx}=u_{yy}$,解得
    \[
        u_{xx}(x,2x)=u_{yy}(x,2x)=-\dfrac{4x}{3},\qquad
        u_{xy}(x,2x)=\dfrac{5x}{3}.
    \]
\end{solution}

\begin{example}
    已知 $f$ 有二阶连续偏导数, 且
    \[
        f(x,x^2)=x^2,\qquad f_x(x,x^2)=kx,\qquad f_{xx}(x,x^2)=2x^2f_{yy}(x,x^2),
    \]
    计算
    \[
        4f_y(x,x^2)-f_{xx}(x,x^2)+x^2f_{yy}(x,x^2).
    \]
\end{example}
\exampleblank

\begin{solution}
    沿 $y=x^2$ 对前两式求导,并使用 $f_{xx}=2x^2f_{yy}$ 消元,得到 $f_y=(2-k)/2$,$x^2f_{yy}=-k/2$.故所求为 $4-3k/2$.
\end{solution}

\begin{example}
    设 $u=f(z)$, 其中 $z=x+y\varphi(z)$, 证明 Lagrange 公式
    \[
        \dfrac{\partial^nu}{\partial y^n}
        =\dfrac{\partial^{n-1}}{\partial x^{n-1}}
        \left[(\varphi(z))^n\dfrac{\partial u}{\partial x}\right].
    \]
\end{example}
\exampleblank

\begin{solution}
    由 $z_y=\varphi(z)/(1-y\varphi'(z))$ 与 $z_x=1/(1-y\varphi'(z))$ 可得 $z_y=\varphi(z)z_x$.对 $n$ 归纳并反复交换 $x,y$ 偏导,即得题给 Lagrange 公式.
\end{solution}

\begin{example}
    设 $u(x,y)$ 有连续二阶偏导数且 $\Delta u=0$, 证明
    \[
        v(x,y)=u\left(\dfrac{x}{x^2+y^2},\dfrac{y}{x^2+y^2}\right)
    \]
    满足 $\Delta v=0$.
\end{example}
\exampleblank

\begin{solution}
    反演映射的微分是比例因子乘正交反射.直接链式求导可得
    $\Delta v=(x^2+y^2)^{-2}(\Delta u)\circ I=0$.
\end{solution}

\begin{example}
    设 $u=f(x,y)$ 有二阶连续偏导数, 证明在极坐标 $x=r\cos\theta,y=r\sin\theta$ 下
    \[
        u_{xx}+u_{yy}=u_{rr}+\dfrac{1}{r}u_r+\dfrac{1}{r^2}u_{\theta\theta}.
    \]
\end{example}
\exampleblank

\begin{solution}
    由 $\partial_x=\cos\theta\,\partial_r-(\sin\theta/r)\partial_\theta$,$\partial_y=\sin\theta\,\partial_r+(\cos\theta/r)\partial_\theta$ 展开,两式相加后交叉项抵消,即得公式.
\end{solution}

\begin{example}
    设 $u=f(x,y,z)$ 有连续二阶偏导数, 证明 Laplace 算子在正交变换下形式不变.
\end{example}
\exampleblank

\begin{solution}
    正交变换矩阵为 $Q$ 时,Hessian 变为 $Q^THQ$;Laplace 算子是 Hessian 的迹,而迹在正交相似变换下不变.
\end{solution}

\begin{example}
    设 $z=f(x,y)$ 在 $\mathbb{R}^2$ 上可微, 若
    \[
        b\dfrac{\partial z}{\partial x}=a\dfrac{\partial z}{\partial y},
        \qquad a\ne0,\ b\ne0,
    \]
    证明 $z$ 是 $ax+by$ 的函数.
\end{example}
\exampleblank

\begin{solution}
    方程说明沿方向 $(b,-a)$ 的方向导数为零,故 $z$ 在直线 $ax+by=C$ 上为常数,于是存在单变量函数 $\Phi$ 使 $z=\Phi(ax+by)$.
\end{solution}

\begin{example}
    设 $f(t)$ 在 $\mathbb{R}$ 上可导, 若对任一调和函数 $u(x,y)$, $F(x,y)=f(u(x,y))$ 仍调和, 证明 $f(t)=at+b$.
\end{example}
\exampleblank

\begin{solution}
    取调和函数 $u=x$,则 $\Delta f(u)=f''(x)=0$,故 $f''\equiv0$,所以 $f(t)=at+b$.
\end{solution}

\begin{example}
    设 $z=f(x,y)$ 满足
    \[
        z_{xx}-z_{yy}+\lambda z_x+\lambda z_y=0,
    \]
    确定 $\alpha,\beta$, 使在变换 $u=ze^{\alpha x+\beta y}$ 下消去一阶偏导项.
\end{example}
\exampleblank

\begin{solution}
    写 $z=u e^{-\alpha x-\beta y}$.一阶项系数分别为 $-2\alpha+\lambda$ 与 $2\beta+\lambda$,故取 $\alpha=\lambda/2,\beta=-\lambda/2$.
\end{solution}

\begin{example}[\markstar]
    \begin{enumerate}
        \item 设 $z$ 满足 $6z_{xx}+z_{xy}-z_{yy}=0$, 确定 $a$, 使在 $u=x-2y,v=x+ay$ 下方程化为 $z_{uv}=0$;
        \item 设 $u$ 满足
              \[
                  Au_{xx}+2Bu_{xy}+Cu_{yy}=0,
                  \qquad AC-B^2<0,\ C\ne0,
              \]
              若在 $\xi=x+\lambda_1y,\eta=x+\lambda_2y$ 下化为 $u_{\xi\eta}=0$, 证明 $\lambda_1,\lambda_2$ 是 $C\lambda^2+2B\lambda+A=0$ 的相异根.
    \end{enumerate}
\end{example}
\exampleblank

\begin{solution}
    (1) 链式展开后纯 $u$ 项自动消失,纯 $v$ 项系数为 $6+a-a^2$;排除退化值 $a=-2$,得 $a=3$.(2) 展开后二阶纯项系数正是 $C\lambda_i^2+2B\lambda_i+A$,故二者为该方程两根;双曲条件保证根相异.
\end{solution}

\begin{example}[\markstar]
    求下列变换的 Jacobi 行列式:
    \begin{enumerate}
        \item $x_1=r\cos\theta,x_2=r\sin\theta$;
        \item $x_1=r\cos\theta_1,x_2=r\sin\theta_1\cos\theta_2,x_3=r\sin\theta_1\sin\theta_2$;
        \item $m$ 维球坐标变换, 并用数学归纳法求 $\dfrac{\partial(x_1,\ldots,x_m)}{\partial(r,\theta_1,\ldots,\theta_{m-1})}$.
    \end{enumerate}
\end{example}
\exampleblank

\begin{solution}
    (1) $J=r$.(2) $J=r^2\sin\theta_1$.(3) 归纳展开最后一列,得到
    \[
        J_m=r^{m-1}\sin^{m-2}\theta_1\sin^{m-3}\theta_2\cdots\sin\theta_{m-2}
    \]
    (符号取决于角变量排列,体积公式取绝对值).
\end{solution}

\begin{example}
    利用 $u=xy,v=x,w=xz-y$ 将 $yz_{yy}+2z_y=x^2$ 化为 $w$ 关于 $u,v$ 的方程.
\end{example}
\exampleblank

\begin{solution}
    由 $y=u/v,z=w/v+u/v^2$,且固定 $x=v$ 时 $\partial_y=v\partial_u$.代入得
    \[
        u w_{uu}+2w_u=v^2-\dfrac{2}{v}.
    \]
\end{solution}

\begin{example}
    利用 $x=u$, $\displaystyle y=\dfrac{u}{1+uv}$, $\displaystyle z=\dfrac{u}{1+uw}$, 将 $x^2z_x+y^2z_y=z^2$ 化为 $w$ 关于 $u,v$ 的方程.
\end{example}
\exampleblank

\begin{solution}
    反解 $v=1/y-1/x,w=1/z-1/x$.向量场 $L=x^2\partial_x+y^2\partial_y$ 满足 $Lu=u^2,Lv=0$;原方程又给出 $Lw=0$,故新方程为 $w_u=0$.
\end{solution}

\begin{proposition}[\markstar]
    设
    \[
        \begin{aligned}
            F(\Delta x,\Delta y)={} & f(x_0+\Delta x,y_0+\Delta y)-f(x_0+\Delta x,y_0) \\
                                    & -f(x_0,y_0+\Delta y)+f(x_0,y_0),
        \end{aligned}
    \]
    在足够光滑条件下, 可由二元中值定理刻画混合偏导, 并得到混合偏导交换的相应充分条件.
\end{proposition}

\begin{example}
    设 $f$ 在 $(x_0,y_0)$ 某邻域存在 $f_{xy},f_{yx}$, 且 $f_{xy}$ 在该点对 $x$ 连续, $f_{yx}$ 在该点关于 $y$ 连续, 证明 $f_{xy}(x_0,y_0)=f_{yx}(x_0,y_0)$.
\end{example}
\exampleblank

\begin{solution}
    对矩形增量 $\Delta_h\Delta_kf$ 分别按两种次序使用一元中值定理,得到两个混合偏导在邻近点的值;令 $h,k\to0$,用题设的单向连续性即得相等.
\end{solution}

\begin{example}
    设 $f_x,f_{xy},f_{yx}$ 在 $(x_0,y_0)$ 某邻域存在, $f_{yx}$ 在该点连续, 证明题设下的混合偏导相等结论.
\end{example}
\exampleblank

\begin{solution}
    把 $f_x(x_0+h,y_0+k)-f_x(x_0+h,y_0)$ 与 $f_x(x_0,y_0+k)-f_x(x_0,y_0)$ 比较,分别用中值定理和 $f_{yx}$ 的连续性,除以 $k$ 后令增量趋零,即得结论.
\end{solution}

\begin{example}
    设 $f_x,f_y$ 在 $(x_0,y_0)$ 某邻域存在, 且在该点可微, 证明 $f_{xy}(x_0,y_0)=f_{yx}(x_0,y_0)$.
\end{example}
\exampleblank

\begin{solution}
    $f_x,f_y$ 在该点可微,故其一阶增量中的交叉系数分别为 $f_{xy},f_{yx}$.而 $f$ 的二阶矩形增量与交换 $h,k$ 对称,比较两种展开便得二者相等.
\end{solution}

\begin{definition}
    若 $f(tx,ty)=t^nf(x,y)$ $(t>0)$, 则称 $f$ 为 $n$ 次齐次函数.
\end{definition}

\begin{proposition}[\markstar]
    设 $f$ 有连续偏导数, 则 $f$ 为 $n$ 次齐次函数当且仅当
    \[
        x\dfrac{\partial f}{\partial x}+y\dfrac{\partial f}{\partial y}=nf(x,y).
    \]
\end{proposition}

\begin{proposition}[\markstar]
    设 $f$ 是二次连续可微的 $n$ 次齐次函数, 则
    \[
        \left(x\dfrac{\partial}{\partial x}+y\dfrac{\partial}{\partial y}\right)^2f=n^2f,
    \]
    且 $f_x,f_y$ 为 $n-1$ 次齐次函数.
\end{proposition}

\begin{proposition}
    若 $f$ 是 $\mathbb{R}^2\setminus\{\bm{0}\}$ 上的连续 $n$ 次齐次函数, 则存在 $C>0$ 使
    \[
        |f(x,y)|\leq C(x^2+y^2)^{\frac{n}{2}}.
    \]
\end{proposition}

\begin{example}
    证明
    \[
        u=\varphi\left(\dfrac{x}{y}\right)+y\psi\left(\dfrac{x}{y}\right)
    \]
    满足方程
    \[
        x^2\dfrac{\partial^2u}{\partial x^2}
        +2xy\dfrac{\partial^2u}{\partial x\partial y}
        +y^2\dfrac{\partial^2u}{\partial y^2}=0,
    \]
    其中 $\varphi,\psi$ 均为二次连续可微函数.
\end{example}
\exampleblank

\begin{solution}
    记 Euler 算子 $E=x\partial_x+y\partial_y$.题给 $u$ 是零次齐次函数与一次齐次函数之和,故 $E(E-1)u=0$;展开即为题给二阶方程.
\end{solution}



\subsection{几何应用}

\begin{proposition}[\markstar]
    设空间曲线 $l$ 在 $(x_0,y_0,z_0)$ 处的切向量为 $\bm\tau=(\tau_x,\tau_y,\tau_z)$, 则切线和法平面分别为
    \[
        \dfrac{x-x_0}{\tau_x}=\dfrac{y-y_0}{\tau_y}=\dfrac{z-z_0}{\tau_z},
    \]
    \[
        \tau_x(x-x_0)+\tau_y(y-y_0)+\tau_z(z-z_0)=0.
    \]
    若曲线由 $F(x,y,z)=G(x,y,z)=0$ 给出, 可取
    \[
        \bm\tau=
        \left(
        \dfrac{\partial(F,G)}{\partial(y,z)},
        \dfrac{\partial(F,G)}{\partial(z,x)},
        \dfrac{\partial(F,G)}{\partial(x,y)}
        \right).
    \]
\end{proposition}

\begin{proposition}[\markstar]
    设空间曲面 $S$ 在 $(x_0,y_0,z_0)$ 处法向量为 $\bm n=(A,B,C)$, 则法线和切平面分别为
    \[
        \dfrac{x-x_0}{A}=\dfrac{y-y_0}{B}=\dfrac{z-z_0}{C},
    \]
    \[
        A(x-x_0)+B(y-y_0)+C(z-z_0)=0.
    \]
    对于 $F(x,y,z)=0$, 可取 $\bm n=\nabla F$; 对参数曲面可取两个切向量的叉积.
\end{proposition}

\begin{example}
    \begin{enumerate}
        \item 求曲线
              \[
                  x=a\sin^2t,\qquad y=b\sin t\cos t,\qquad z=c\cos^2t
              \]
              在 $t=\dfrac{\pi}{4}$ 处的切线方程和法平面方程;
        \item 求曲线
              \[
                  x^2+y^2+z^2=4a^2,\qquad x^2+y^2=2ax
              \]
              在点 $(a,a,\sqrt2a)$ 处的切线方程与法平面方程.
    \end{enumerate}
\end{example}
\exampleblank

\begin{solution}
    (1) 点为 $(a/2,b/2,c/2)$,切向量为 $(a,0,-c)$.切线据此写出,法平面为 $a(x-a/2)-c(z-c/2)=0$.(2) 两曲面法向量叉乘得切向量 $(-\sqrt{2},0,1)$,故切线为该点加其参数倍,法平面为 $-\sqrt{2}(x-a)+z-\sqrt{2}a=0$.
\end{solution}

\begin{example}
    \begin{enumerate}
        \item 求曲面 $x=u\cos v,y=u\sin v,z=v$ 在 $(\sqrt{2},\sqrt{2},\dfrac{\pi}{4})$ 处的切平面;
        \item 求 $x^2+2y^2+3z^2=21$ 在 $(1,-2,2)$ 处的法线;
        \item 求该曲面平行于平面 $x+4y+6z=0$ 的切平面.
    \end{enumerate}
\end{example}
\exampleblank

\begin{solution}
    (1) 切平面为 $x-y+2\sqrt{2}(z-\pi/4)=0$.(2) 法线为 $(x,y,z)=(1,-2,2)+t(1,-4,6)$.(3) 切点为 $(1,2,2)$ 与 $(-1,-2,-2)$,相应平面为 $x+4y+6z=21$ 与 $x+4y+6z=-21$.
\end{solution}

\begin{example}
    \begin{enumerate}
        \item 设 $F(u,v)$ 可微, 证明曲面 $\displaystyle F\left(\dfrac{x-a}{z-c},\dfrac{y-b}{z-c}\right)=0$ 的所有切平面过一定点;
        \item 设 $F(u,v)$ 可微, 给定曲面 $F(nx-ly,ny-mz)=0$, 证明其所有切平面都与某定直线平行.
    \end{enumerate}
\end{example}
\exampleblank

\begin{solution}
    (1) 该式只依赖从 $(a,b,c)$ 出发的射线方向,Euler 关系给出切平面均过 $(a,b,c)$.(2) 取方向向量 $(lm,mn,n^2)$,它同时使两个线性自变量的增量为零,故所有切平面都平行于这条定直线.
\end{solution}

\begin{example}
    设 $F(x,y,z)$ 是可微的 $n$ 次齐次函数, 证明曲面 $F(x,y,z)=1$ 在 $(x_0,y_0,z_0)$ 处的切平面为
    \[
        xF_x(x_0,y_0,z_0)+yF_y(x_0,y_0,z_0)+zF_z(x_0,y_0,z_0)=n.
    \]
\end{example}
\exampleblank

\begin{solution}
    切平面为 $F_x(X_0)x+F_y(X_0)y+F_z(X_0)z=F_xx_0+F_yy_0+F_zz_0$;Euler 齐次公式把右端化为 $nF(X_0)=n$.
\end{solution}



\subsection{Taylor 公式}

\begin{theorem}[\markstar\ 带 Lagrange 余项的 Taylor 公式]
    设 $D$ 是平面凸域, $f\in C^{(n+1)}(D)$, $(x_0,y_0)\in D$, 则对 $(x_0+h,y_0+k)\in D$,
    \[
        \begin{aligned}
            f(x_0+h,y_0+k)
            ={} & \sum\limits_{m=0}^n\dfrac{1}{m!}
            \left(h\dfrac{\partial}{\partial x}+k\dfrac{\partial}{\partial y}\right)^m f(x_0,y_0) \\
                & +\dfrac{1}{(n+1)!}
            \left(h\dfrac{\partial}{\partial x}+k\dfrac{\partial}{\partial y}\right)^{n+1}
            f(x_0+\theta h,y_0+\theta k),
        \end{aligned}
    \]
    其中 $0<\theta<1$, 且
    \[
        \left(h\dfrac{\partial}{\partial x}+k\dfrac{\partial}{\partial y}\right)^mf
        =\sum\limits_{r=0}^mC_m^r h^{m-r}k^r
        \dfrac{\partial^mf}{\partial x^{m-r}\partial y^r}.
    \]
\end{theorem}

\begin{theorem}[\markstar\ 带 Peano 余项的 Taylor 公式]
    若 $f\in C^{(n)}(D)$, 则
    \[
        f(x_0+h,y_0+k)
        =\sum\limits_{m=0}^n\dfrac{1}{m!}
        \left(h\dfrac{\partial}{\partial x}+k\dfrac{\partial}{\partial y}\right)^m f(x_0,y_0)
        +o(\rho^n),
    \]
    其中 $\rho=\sqrt{h^2+k^2}\to0$.
\end{theorem}

\begin{example}[\markstar]
    若 $f$ 具有 $n+1$ 阶连续偏导数, 且
    \[
        f(x,y)=\sum\limits_{i+j=0}^nA_{ij}(x-x_0)^i(y-y_0)^j+o(\rho^n),
    \]
    证明 Taylor 展开的系数唯一, 且
    \[
        A_{ij}=\dfrac{1}{i!j!}
        \dfrac{\partial^{i+j}f}{\partial x^i\partial y^j}(x_0,y_0).
    \]
\end{example}
\exampleblank

\begin{solution}
    逐次对展开式作不超过 $n$ 阶偏导并令 $(x,y)\to(x_0,y_0)$;最低未消失项唯一给出相应系数.归纳得到 $i!j!A_{ij}=\partial_x^i\partial_y^jf(x_0,y_0)$.
\end{solution}

\begin{example}
    \begin{enumerate}
        \item 求 $f(x,y)=e^x\cos y$ 在 $(0,0)$ 处带 Peano 余项的 Taylor 展开至四阶;
        \item 求 $f(x,y)=y^2$ 在 $(1,1)$ 处带 Peano 余项的 Taylor 展开至二阶.
    \end{enumerate}
\end{example}
\exampleblank

\begin{solution}
    (1)
    \[
        e^x\cos y=1+x+\dfrac{x^2-y^2}{2}+\dfrac{x^3-3xy^2}{6}
        +\dfrac{x^4-6x^2y^2+y^4}{24}+o(r^4).
    \]
    (2) 令 $h=x-1,k=y-1$,则 $y^2=1+2k+k^2$.
\end{solution}

\begin{example}
    求 $f(x,y)=2x^2-xy+y^2$ 在 $(1,-2)$ 处的 Taylor 展开式.
\end{example}
\exampleblank

\begin{solution}
    令 $h=x-1,k=y+2$,精确展开为
    \[
        f=8+6h-5k+2h^2-hk+k^2.
    \]
\end{solution}

\begin{example}
    设 $D\subseteq\mathbb{R}^n$ 为凸有界闭区域, $f(P)$ 在 $D$ 上有连续一阶偏导数, 证明存在 $L>0$, 对任意 $P_1,P_2\in D$,
    \[
        |f(P_2)-f(P_1)|\leq L|P_2-P_1|.
    \]
\end{example}
\exampleblank

\begin{solution}
    $\nabla f$ 在紧集 $D$ 上有界,记其上界为 $L$.沿 $P_1P_2$ 使用一元中值公式即得题式.
\end{solution}

\begin{example}
    设 $F(x,y,z)$ 在 $\mathbb{R}^3$ 中有连续一阶偏导数, 且
    \[
        yF_x-xF_y+F_z\geq\alpha>0.
    \]
    当 $(x,y,z)$ 沿曲线 $\Gamma:x=-\cos t,y=\sin t,z=t$ $(t\geq0)$ 趋向无穷远时, 证明 $F(x,y,z)\to+\infty$.
\end{example}
\exampleblank

\begin{solution}
    沿曲线有
    \[
        \dfrac{\mathrm d}{\mathrm dt}F(-\cos t,\sin t,t)=F_x\sin t+F_y\cos t+F_z
        =yF_x-xF_y+F_z\geq\alpha.
    \]
    积分得 $F\geq F(0)+\alpha t\to+\infty$.
\end{solution}

\begin{example}
    设 $D$ 是 $\mathbb{R}^2$ 中凸区域, $f\in C^{(2)}(D)$, 讨论下列关于凸函数的一阶,二阶判别条件之间的等价关系.
\end{example}
\exampleblank

\begin{solution}
    对 $C^2$ 函数,凸性等价于一阶支撑不等式 $f(y)\geq f(x)+\nabla f(x)\cdot(y-x)$,也等价于 Hessian 半正定.两者分别沿线段对一元函数应用凸性判别;严格正定给出严格凸,但严格凸不要求处处正定.
\end{solution}


\newpage

\section{多元函数微分学的应用}


\subsection{隐函数存在定理}

\begin{theorem}[\markstar]
    设方程组
    \[
        f(x,y,u,v)=0,\qquad g(x,y,u,v)=0
    \]
    确定 $u=\varphi(x,y),v=\psi(x,y)$, 且
    \[
        J=\dfrac{\partial(f,g)}{\partial(u,v)}\ne0.
    \]
    则隐函数的一阶偏导可由对方程组分别求导后解线性方程得到; 等价地可用 Cramer 法则写成相应 Jacobi 行列式之比.
\end{theorem}

\begin{remark}
    隐函数求导公式不建议机械记忆. 对给定方程组直接两边求导再解方程通常更稳妥. 一般情况下, 有几个独立方程就对应几个因变量.
\end{remark}



\methodsection{(A) 计算}

\begin{example}
    \begin{enumerate}
        \item 设 $z=z(x,y)$ 由 $z^3-3xyz=a^3$ 确定, 求 $z_{xy}$ $(z^2\ne xy)$;
        \item 设 $z=z(x,y)$ 由 $z^x=y^z$ 确定, 求 $z_{xx}$.
    \end{enumerate}
\end{example}
\exampleblank

\begin{solution}
    (1) 令 $F=z^3-3xyz-a^3$,隐式求导得
    \[
        z_{xy}=-\dfrac{z(x^2y^2+2xyz^2-z^4)}{(z^2-xy)^3}.
    \]
    (2) 令 $F=x\ln z-z\ln y$,$D=F_z=x/z-\ln y$,则
    \[
        z_x=-\dfrac{\ln z}{D},\qquad
        z_{xx}=-\dfrac{2z_x/z-xz_x^2/z^2}{D}.
    \]
\end{solution}

\begin{example}
    设 $x^2-xy+2y^2+x-y-1=0$, 求 $x=0,y=1$ 时 $y',y'',y'''$.
\end{example}
\exampleblank

\begin{solution}
    把 $y=1+y'(0)x+y''(0)x^2/2+y'''(0)x^3/6+o(x^3)$ 代回方程并比较系数,得
    \[
        y'(0)=0,\qquad y''(0)=-\dfrac{2}{3},\qquad y'''(0)=-\dfrac{2}{3}.
    \]
\end{solution}

\begin{example}
    \begin{enumerate}
        \item 设 $f$ 可微, $z$ 由 $f(x+y+z,x^2+y^2+z^2)=0$ 确定, 求 $z_x$;
        \item 设 $f(u)$ 可微, 方程 $y=f(x^2+y^2)+f(x+y)$ 确定 $y=g(x)$, 已知 $g(0)=2$, $f'(2)=\dfrac{1}{2}$, $f'(4)=1$, 求 $g'(0)$.
    \end{enumerate}
\end{example}
\exampleblank

\begin{solution}
    (1) 记 $f_1,f_2$ 为两个偏导,则
    \[
        z_x=-\dfrac{f_1+2xf_2}{f_1+2zf_2}.
    \]
    (2) 对方程求导并代入给定数据,得到 $g'(0)=-1/7$.
\end{solution}

\begin{example}
    \begin{enumerate}
        \item 设 $z=z(x,y)$ 由 $\displaystyle x^z=\ln\dfrac{z}{y}+1$ 确定, 求 $\mathrm{d}^2z$;
        \item 设 $F(u,v)$ 二次可微, $z$ 由 $\displaystyle F\left(\dfrac{x}{z},\dfrac{y}{z}\right)=0$ 确定, 求 $\mathrm{d}^2z$.
    \end{enumerate}
\end{example}
\exampleblank

\begin{solution}
    两问统一使用隐函数二阶微分公式.若方程写成 $G(x,y,z)=0$,则
    \[
        \mathrm dz=-\dfrac{G_x\mathrm dx+G_y\mathrm dy}{G_z},
    \]
    \[
        \mathrm d^2z=-\dfrac{G_{xx}\mathrm dx^2+2G_{xy}\mathrm dx\mathrm dy+G_{yy}\mathrm dy^2
            +2G_{xz}\mathrm dx\mathrm dz+2G_{yz}\mathrm dy\mathrm dz+G_{zz}\mathrm dz^2}{G_z}.
    \]
    分别取 $G=x^z-\ln(z/y)-1$ 与 $G=F(x/z,y/z)$ 即得所求.
\end{solution}

\begin{example}
    设 $f(x,y,z)$ 是 $n$ 次齐次可微函数, 若 $f(x,y,z)=0$ 隐含 $z=\varphi(x,y)$, 证明 $\varphi$ 是一次齐次函数.
\end{example}
\exampleblank

\begin{solution}
    由齐次性,若 $f(x,y,z)=0$,则 $f(tx,ty,tz)=0$.局部解的唯一性给出 $\varphi(tx,ty)=t\varphi(x,y)$,故为一次齐次.
\end{solution}

\begin{example}
    \begin{enumerate}
        \item 设 $f(u,v,w),g(u,v,w),h(u,v,w)$ 可微, $u=u(x,y,z)$ 由方程组
              \[
                  \begin{cases}
                      x=f(u,v,w), \\
                      y=g(u,v,w), \\
                      z=h(u,v,w)
                  \end{cases}
              \]
              所确定, 求 $\dfrac{\partial u}{\partial x},\dfrac{\partial u}{\partial y},\dfrac{\partial u}{\partial z}$;
        \item 设 $f(x,y,z,t),g(y,z,t),h(z,t)$ 可微, $u(x,y)$ 由
              \[
                  \begin{cases}
                      u=f(x,y,z,t), \\
                      g(y,z,t)=0,   \\
                      h(z,t)=0
                  \end{cases}
              \]
              所确定, 若 $\dfrac{\partial(g,h)}{\partial(z,t)}\ne0$, 求 $\dfrac{\partial u}{\partial y}$;
        \item 设 $u=f(x,y,z)$, $g(x^2,e^y,z)=0$, $y=\sin x$, 且 $f,g$ 都有一阶连续偏导数, 求 $\dfrac{\partial u}{\partial x}$.
    \end{enumerate}
\end{example}
\exampleblank

\begin{solution}
    (1) 记 $J=\partial(f,g,h)/\partial(u,v,w)$,则
    $u_x=\partial(g,h)/\partial(v,w)/J$,$u_y=-\partial(f,h)/\partial(v,w)/J$,$u_z=\partial(f,g)/\partial(v,w)/J$.(2) 记 $J=\partial(g,h)/\partial(z,t)$,则
    $u_y=f_y+g_y(f_th_z-f_zh_t)/J$.(3)
    \[
        u_x=f_x+f_y\cos x-f_z\dfrac{2xg_1+e^y\cos x\,g_2}{g_3}.
    \]
\end{solution}



\methodsection{(B) 理论证明 *}

\begin{proposition}[\markstar\ 整体解]
    设
    \[
        D=\{(x,y)\mid a<x<b,\ -\infty<y<+\infty\},
    \]
    $F$ 在 $D$ 上连续, 且 $F_y\geq m>0$. 方程 $F(x,y)=0$ 在 $(a,b)$ 上存在唯一连续解 $y=y(x)$.
\end{proposition}

\begin{example}
    设 $D=\{(x,y)\mid0<x<b,-\infty<y<+\infty\}$,
    \[
        f(x,y)=1-e^{-x}+\dfrac{y^3}{e^y},
    \]
    证明存在 $(0,+\infty)$ 上的 $C^1$ 函数 $y(x)$ 满足 $f(x,y)=0$.
\end{example}
\exampleblank

\begin{solution}
    对固定 $x>0$,方程等价于 $y^3e^{-y}=e^{-x}-1<0$.左端在 $(-\infty,0)$ 严格递增并从 $-\infty$ 到 $0$,故有唯一负根;其对 $y$ 的导数不为零,隐函数定理给出全区间上的 $C^1$ 解.
\end{solution}

\begin{example}
    设 $f(x,y)=y\cos^2x-\sin^2x-e^{-y}$, 证明不存在 $(-\infty,+\infty)$ 上的隐函数 $y=y(x)$ 满足 $f(x,y)=0$.
\end{example}
\exampleblank

\begin{solution}
    取 $x=\pi/2$,方程左端为 $-1-e^{-y}<0$,不可能为零,故不存在定义于全实轴的隐函数.
\end{solution}

\begin{example}
    设 $F$ 在矩形 $D=\lbrack a,b\rbrack\times\lbrack c,d\rbrack$ 上满足题设偏导与符号条件, 证明 $F(x,y)=0$ 在 $\lbrack a,b\rbrack$ 上存在唯一可微解 $y=y(x)$.
\end{example}
\exampleblank

\begin{solution}
    题面没有列出所谓偏导与符号条件,无法验证存在性,唯一性或可微性.

    通常需要 $F(x,c)F(x,d)<0$,并要求 $F_y$ 保持非零同号.在这些条件下,由介值定理,单调性和隐函数定理即得结论.
\end{solution}

\begin{example}
    设 $\varphi$ 在 $\lbrack-a,a\rbrack$ 上可微, $\varphi(0)=0$, $|\varphi'(y)|\leq k<1$, 证明存在 $\delta>0$, 使方程 $x=y+\varphi(y)$ 在 $\lbrack-\delta,\delta\rbrack$ 上有唯一可微解 $y=y(x)$, 且 $y(0)=0$.
\end{example}
\exampleblank

\begin{solution}
    $H(y)=y+\varphi(y)$ 满足 $H'(y)\geq1-k>0$,故严格递增;在零点附近其像包含某区间 $[-\delta,\delta]$.逆函数定理给出唯一可微逆函数,且 $y(0)=0$.
\end{solution}

\begin{example}
    设 $f(x)$ 在 $\lbrack a,b\rbrack$ 上可微, $R(x,y)$ 在 $\{(x,y)\mid a\leq y\leq x\leq b\}$ 上连续, 证明积分方程
    \[
        \varphi(x)=f(x)+\lambda\int_a^xR(x,y)\varphi(y)\,\mathrm{d}y
    \]
    在 $\lbrack a,b\rbrack$ 上有唯一连续解.
\end{example}
\exampleblank

\begin{solution}
    以 $\varphi_0=f$ 开始作 Picard 迭代.若 $|R|\leq M$,第 $n$ 次差由 $(|\lambda|M)^n(x-a)^n/n!$ 控制,故级数一致收敛到连续解;同一估计用于两解之差即得唯一性.
\end{solution}

\begin{example}
    设 $f,g$ 在区域 $D\subseteq\mathbb{R}^2$ 上有连续偏导数, 且
    \[
        \dfrac{\partial(f,g)}{\partial(x,y)}\ne0.
    \]
    若 $D'\subseteq D$ 是有界闭区域, 证明 $D'$ 内同时满足 $f=0,g=0$ 的点只有有限个.
\end{example}
\exampleblank

\begin{solution}
    Jacobian 非零使每个公共零点由逆函数定理成为孤立点.若有无穷多个,紧集 $D'$ 中必有聚点,与孤立性矛盾.
\end{solution}

\begin{example}
    设 $f$ 在 $\mathbb{R}^2$ 上有连续二阶偏导数, $X_0$ 为驻点, 且 Hesse 行列式在 $X_0$ 处非零, 证明存在 $X_0$ 的邻域, 其中除 $X_0$ 外无其他驻点.
\end{example}
\exampleblank

\begin{solution}
    驻点是映射 $T=(f_x,f_y)$ 的零点;其 Jacobian 正是非零 Hesse 行列式.逆函数定理说明 $T$ 在邻域内一一,故零点唯一.
\end{solution}

\begin{example}
    设 $T:\mathbb{R}^2\to\mathbb{R}^2$ 为连续可微映射, Jacobi 行列式处处非零, 且任一有界闭集 $K$ 的逆像 $T^{-1}(K)$ 仍为有界闭集, 证明 $T$ 为满射.
\end{example}
\exampleblank

\begin{solution}
    局部逆定理说明 $T(\mathbb R^2)$ 为开集.题设的 proper 性说明收敛像列的原像有收敛子列,从而像也是闭集.它非空,而 $\mathbb R^2$ 连通,故像为全平面.
\end{solution}

\begin{example}
    设 $f\in C^1(\mathbb{R})$, 映射 $T:(x,y)\mapsto(u,v)=(f(x),-y+xf(x))$, 若 $f'(x_0)\ne0$, 证明 $T$ 在 $(x_0,y_0)$ 某邻域可逆, 并写出局部逆映射.
\end{example}
\exampleblank

\begin{solution}
    Jacobian 为 $-f'(x)$,在给定点非零.先局部解 $x=f^{-1}(u)$,再由 $v=-y+xu$ 得
    \[
        x=f^{-1}(u),\qquad y=u f^{-1}(u)-v.
    \]
\end{solution}

\begin{example}
    设 $k\ne0$, $\varphi(y)$ 为周期 $T$ 的可微函数, 且 $|\varphi'(y)|\leq|k|$. 若 $x=ky+\varphi(y)$ 确定隐函数 $y=y(x)$, 证明存在周期为 $|k|T$ 的函数 $\psi(x)$, 使
    \[
        y=\dfrac{x}{k}+\psi(x).
    \]
\end{example}
\exampleblank

\begin{solution}
    令 $H(y)=ky+\varphi(y)$.由隐函数存在假设及 $H(y+T)=H(y)+kT$,其逆满足 $y(x+kT)=y(x)+T$.于是 $\psi(x)=y(x)-x/k$ 具有周期 $|k|T$(改变周期方向不影响绝对值).
\end{solution}

\begin{example}
    设 $f$ 有二阶连续偏导数, $T:(x,y)\mapsto(u,v)=(f_x,f_y)$, 若
    \[
        f_{xx}f_{yy}-f_{xy}^2\ne0,
    \]
    证明局部存在逆映射 $T^{-1}:(u,v)\mapsto(x,y)$, 且存在函数 $g$ 满足 $x=g_u,y=g_v$.
\end{example}
\exampleblank

\begin{solution}
    逆函数定理给出 $(x,y)=(x(u,v),y(u,v))$.定义 Legendre 函数
    $g(u,v)=ux+vy-f(x,y)$,微分中利用 $u=f_x,v=f_y$ 得 $\mathrm dg=x\,\mathrm du+y\,\mathrm dv$,故 $x=g_u,y=g_v$.
\end{solution}



\methodsection{(C) 秩定理 *}

\begin{theorem}[\markstar\ 秩定理]
    设 $A\subseteq\mathbb{R}^m$, $B\subseteq\mathbb{R}^n$ 为开集, $F:A\to B$ 是 $C^k$ 映射, 且 $F$ 在 $A$ 上秩恒为 $r$. 对 $a\in A$, $b=F(a)$, 存在 $a,b$ 的开邻域和 $C^k$ 微分同胚 $u,v$, 使
    \[
        v\circ F\circ u^{-1}(x_1,\ldots,x_m)
        =(x_1,\ldots,x_r,0,\ldots,0).
    \]
\end{theorem}

\begin{example}
    讨论秩定理在满秩嵌入情形下的特殊形式, 即将 $C^k$ 映射局部化为
    \[
        (x_1,\ldots,x_m)\mapsto(x_1,\ldots,x_m,0,\ldots,0).
    \]
\end{example}
\exampleblank

\begin{solution}
    满秩嵌入时选取像中 $m$ 个独立坐标与定义域坐标组成局部坐标,逆函数定理把映射化为标准包含 $x\mapsto(x,0)$.
\end{solution}

\begin{example}
    讨论秩定理在满秩投影情形下的特殊形式, 即局部坐标变换后
    \[
        (x_1,\ldots,x_m)\mapsto(x_1,\ldots,x_n).
    \]
\end{example}
\exampleblank

\begin{solution}
    满秩投影时选取 Jacobian 的一个非零 $n$ 阶子式,把相应输出坐标与剩余输入坐标合成局部坐标,逆函数定理把映射化为标准投影.
\end{solution}

\begin{remark}
    可回忆高等代数中的相抵标准型, 从坐标变换和秩不变性的角度比较两者的形式联系.
\end{remark}



\methodsection{(D) 函数相关性 *}

\begin{definition}
    设 $y_i=y_i(x_1,\ldots,x_n)$ $(i=1,\ldots,m)$ 在 $x_0$ 的邻域有定义. 若存在在 $y_0$ 任意邻域内不恒为零的函数 $F(y_1,\ldots,y_m)$, 使在 $x_0$ 某邻域内恒有
    \[
        F(y_1(x),\ldots,y_m(x))\equiv0,
    \]
    则称 $y_1,\ldots,y_m$ 在 $x_0$ 处函数相关, 否则称为函数无关.
\end{definition}

\begin{theorem}[\markstar]
    设 $y_i=y_i(x_1,\ldots,x_n)$ 在 $x_0$ 某邻域有连续一阶偏导数, Jacobi 矩阵在 $x_0$ 附近秩为 $r>0$, 且某个 $r$ 阶子式非零. 则其中对应的 $r$ 个函数在 $x_0$ 处函数无关; 当 $m>r$ 时, 全部 $m$ 个函数在 $x_0$ 处函数相关.
\end{theorem}

\begin{example}
    判别下列函数之间是否函数相关:
    \begin{enumerate}
        \item $u=x^2-y^2$, $v=2xy$;
        \item $u=x+y+z$, $v=xy+yz+zx$, $w=x^2+y^2+z^2$;
        \item $u=\ln x+\ln y$, $v=\cos(xy)$.
    \end{enumerate}
\end{example}
\exampleblank

\begin{solution}
    (1) Jacobian $\partial(u,v)/\partial(x,y)=4(x^2+y^2)$,除原点外独立,故不相关.(2) 有 $w=u^2-2v$,相关.(3) 因 $xy=e^u$,有 $v=\cos(e^u)$,相关.
\end{solution}

\begin{example}
    设 $u=f(x,y)$ 在 $\mathbb{R}^2$ 上可微:
    \begin{enumerate}
        \item 求 $u$ 只是 $x^2+y^2$ 的函数的条件;
        \item 求 $u$ 只是 $\dfrac{y}{x}$ 的函数的条件.
    \end{enumerate}
\end{example}
\exampleblank

\begin{solution}
    (1) 条件为 $yu_x-xu_y=0$;(2) 条件为 $xu_x+yu_y=0$.在相应连通域内,它们分别说明沿圆周或射线方向不变,因而也是充分条件.
\end{solution}

\begin{example}
    设 $f\in C(\mathbb{R})$, 若
    \[
        \int_a^bf(x)\,\mathrm{d}x
    \]
    只依赖于 $\dfrac{b}{a}$, 证明 $f(x)=\dfrac{k}{x}$.
\end{example}
\exampleblank

\begin{solution}
    设积分为 $H(b/a)$.令 $b=ta$ 并对 $a$ 求导,得 $t f(at)-f(a)=0$.因此 $xf(x)$ 为常数,即 $f(x)=k/x$(题意需限于不穿过零点的正区间).
\end{solution}



\subsection{极值}


\methodsection{(A) 自由极值}

设 $P_0$ 为 $f$ 的驻点, Hesse 矩阵为 $Q$. 二阶 Taylor 公式给出
\[
    f(P_0+\Delta)=f(P_0)+\dfrac{1}{2}\Delta\bm x\,Q\,\Delta\bm x^{\mathsf T}+o(r^2).
\]

\begin{theorem}[\markstar]
    设 $f$ 在驻点 $P_0$ 某邻域有二阶连续偏导数:
    \begin{enumerate}
        \item $Q$ 正定时 $P_0$ 为极小值点;
        \item $Q$ 负定时 $P_0$ 为极大值点;
        \item $Q$ 不定时 $P_0$ 不是极值点;
        \item $Q$ 半正定或半负定时需进一步判别.
    \end{enumerate}
\end{theorem}

\begin{remark}
    开集上的极值点只可能来自驻点或偏导不存在的点. 对闭区域最值问题还需考察边界; 对无界区域常先排除明显不可能取得最值的无界部分; 对有界开区域有时先连续延拓到闭包再处理.
\end{remark}

\begin{example}
    求 $f(x,y)=x^4+y^4-2x^2+4xy-2y^2$ 的极值.
\end{example}
\exampleblank

\begin{solution}
    驻点为 $(0,0),(\sqrt{2},-\sqrt{2}),(-\sqrt{2},\sqrt{2})$.原点沿不同方向一正一负,是鞍点;后两点为全局最小点,最小值 $-8$.因函数趋于 $+\infty$,无最大值.
\end{solution}

\begin{example}
    求 $f(x,y)=x^3-4x^2+2xy-y^2$ 在
    \[
        D=\{(x,y)\mid -1\leq x\leq4,\ -1\leq y\leq1\}
    \]
    上的最大值.
\end{example}
\exampleblank

\begin{solution}
    检查内部驻点及四条边界的一元极值,最大值为
    \[
        7,
    \]
    在 $(4,1)$ 处取得.
\end{solution}

\begin{example}
    确定 $f(x,y)=4x+xy^2+y^2$ 在 $D=\{x^2+y^2\leq1\}$ 上的最大值和最小值.
\end{example}
\exampleblank

\begin{solution}
    因 $x+1\geq0$,固定 $x$ 时最大值取 $y^2=1-x^2$,最小值取 $y=0$.化为一元函数后得最大值 $4$ 于 $(1,0)$,最小值 $-4$ 于 $(-1,0)$.
\end{solution}

\begin{example}
    求下列方程确定的 $z=z(x,y)$ 的极值:
    \begin{enumerate}
        \item $2x^2+y^2+z^3+2xy-2x-7y-4z+4=0$;
        \item $(x+y)^2+(y+z)^2+(z+x)^2=3$.
    \end{enumerate}
\end{example}
\exampleblank

\begin{solution}
    (1) 驻点须有 $x=-5/2,y=6$,而 $z$ 是 $z^3-4z-29/2=0$ 的唯一实根;该支在此取局部最大值.(2) 有 $x=y=-z/3$,故 $z=\pm3/2$;最大值为 $3/2$,最小值为 $-3/2$.
\end{solution}

\begin{example}
    求曲面 $z=xy-1$ 上与原点最近的点.
\end{example}
\exampleblank

\begin{solution}
    最小化 $x^2+y^2+(xy-1)^2$.驻点方程只有 $x=y=0$,且函数趋于无穷,故最近点为 $(0,0,-1)$,距离为 $1$.
\end{solution}

\begin{example}
    设三角形 $ABC$ 顶点为 $A=(a_1,b_1)$, $B=(a_2,b_2)$, $C=(a_3,b_3)$, 求平面上一点 $P=(x,y)$, 使 $PA^2+PB^2+PC^2$ 最小.
\end{example}
\exampleblank

\begin{solution}
    平方距离和等于 $3|P-G|^2$ 加常数,其中
    $G=((a_1+a_2+a_3)/3,(b_1+b_2+b_3)/3)$ 为重心,故唯一最小点是 $G$.
\end{solution}

\begin{example}
    证明:
    \begin{enumerate}
        \item $\sin x\sin y\sin(x+y)\leq\dfrac{3\sqrt3}{8}$, $0<x,y<\pi$;
        \item $e^y+x\ln x-x-xy\geq0$, $x\geq1,y\geq0$.
    \end{enumerate}
\end{example}
\exampleblank

\begin{solution}
    (1) 若 $x+y\geq\pi$ 左端非正;否则对对数求驻点,唯一内部最大点为 $x=y=\pi/3$,值得 $3\sqrt{3}/8$.(2) 固定 $x$,关于 $y$ 的凸函数在 $y=\ln x$ 取最小值 $0$.
\end{solution}

\begin{example}
    设 $D\subseteq\mathbb{R}^2$ 为有界区域, $f\in C(\overline D)$ 且在 $D$ 内可微, 若 $f\equiv C$ 在 $\partial D$ 上, 证明 $f$ 在 $D$ 内有驻点.
\end{example}
\exampleblank

\begin{solution}
    若最大值或最小值在内部取得,Fermat 定理直接给出驻点;若二者都只在边界且边界值同为 $C$,则非常值函数必在内部取得严格极值.常值情形任一点都是驻点.
\end{solution}

\begin{example}
    设 $f$ 在 $\mathbb{R}^2$ 上有连续三阶偏导数, 且
    \[
        f_{xx}+f_{yy}>0,
    \]
    证明 $f$ 在 $\mathbb{R}^2$ 内没有极大值点.
\end{example}
\exampleblank

\begin{solution}
    内点极大值要求 Hessian 半负定,因而 $f_{xx}+f_{yy}\leq0$,与题设矛盾.
\end{solution}

\begin{example}
    设 $f$ 在 $X_0=(x_0,y_0)$ 某邻域有连续二阶偏导数, 若 $f$ 在 $X_0$ 达到极大值, 证明对任意 $g,h$,
    \[
        \left(f_{xx}+f_{yy}+gf_x+hf_y\right)_{(x_0,y_0)}\leq0.
    \]
\end{example}
\exampleblank

\begin{solution}
    极大值点处 $f_x=f_y=0$,且 Hessian 半负定,所以 $f_{xx}+f_{yy}\leq0$;含 $g,h$ 的项为零,结论成立.
\end{solution}

\begin{example}
    设 $D=\{(x,y)\mid0<x<1,0<y<1\}$,
    \[
        f(x,y)=ax^2+2bxy+cy^2+dx+ey.
    \]
    若 $f\leq0$ 在 $\partial D$ 上成立, 证明 $f\leq0$ 在 $D$ 内成立.
\end{example}
\exampleblank

\begin{solution}
    把二次多项式减去其在四个顶点的双线性插值,或逐次对 $x,y$ 配方,可将内部值写成边界值的凸组合加上不正的二次余项,从而 $f\leq0$.
\end{solution}



\methodsection{(B) 条件极值}

由 Lagrange 乘数法可将条件极值问题转化为自由极值问题.

\begin{example}
    求 $f(x,y,z)=xyz$ 在
    \[
        x^2+y^2+z^2=1,
        \qquad x+y+z=0
    \]
    下的极值.
\end{example}
\exampleblank

\begin{solution}
    Lagrange 方程给出非零极值点是 $(a,a,-2a)$ 的排列,$6a^2=1$.故
    \[
        \max xyz=\dfrac{1}{3\sqrt{6}},\qquad
        \min xyz=-\dfrac{1}{3\sqrt{6}}.
    \]
\end{solution}

\begin{example}
    证明周长一定的三角形中, 等边三角形面积最大.
\end{example}
\exampleblank

\begin{solution}
    由 Heron 公式,在半周长 $s$ 固定时面积平方为 $s(s-a)(s-b)(s-c)$.三个正数 $(s-a),(s-b),(s-c)$ 和固定,AM--GM 表明相等时乘积最大,即三角形等边.
\end{solution}

\begin{example}
    求 $f(x,y,z)=ax+by+cz$ 在 $x^2+y^2+z^2=1$ 下的最大值与最小值.
\end{example}
\exampleblank

\begin{solution}
    Cauchy--Schwarz 给出 $|ax+by+cz|\leq\sqrt{a^2+b^2+c^2}$,等号在单位向量与 $(a,b,c)$ 同向或反向时取得.
\end{solution}

\begin{example}
    求 $u=x+y+z$ 在
    \[
        \dfrac{1}{x}+\dfrac{4}{y}+\dfrac{9}{z}=1,
        \qquad x,y,z>0
    \]
    下的极值.
\end{example}
\exampleblank

\begin{solution}
    Cauchy--Schwarz 给出
    \[
        (x+y+z)\left(\dfrac{1}{x}+\dfrac{4}{y}+\dfrac{9}{z}\right)\geq(1+2+3)^2=36.
    \]
    故最小值为 $36$,在 $(6,12,18)$ 取得;无最大值.
\end{solution}

\begin{example}[\markstar]
    设 $a>0,c>0,ac-b^2>0$, 方程 $ax^2+2bxy+cy^2=1$ 表示椭圆. 利用 Lagrange 乘数法证明其面积为
    \[
        \dfrac{\pi}{\sqrt{ac-b^2}}.
    \]
\end{example}
\exampleblank

\begin{solution}
    正定二次型可经线性变换化为 $u^2+v^2=1$,变换行列式绝对值为 $\sqrt{ac-b^2}$.单位圆面积除以该值即得 $\pi/\sqrt{ac-b^2}$.
\end{solution}

\begin{example}
    证明:
    \begin{enumerate}
        \item 对 $x,y,z>0$,
              \[
                  xy^2z^3\leq108\left(\dfrac{x+y+z}{6}\right)^6;
              \]
        \item 对 $x,y>0$, $n\in\mathbb{Z}^+$,
              \[
                  \dfrac{x^n+y^n}{2}\geq\left(\dfrac{x+y}{2}\right)^n.
              \]
    \end{enumerate}
\end{example}
\exampleblank

\begin{solution}
    (1) 对六个数 $x/1,y/2,y/2,z/3,z/3,z/3$ 用 AM--GM,整理即得.(2) 函数 $t^n$ 在正半轴为凸函数,Jensen 不等式给出结论.
\end{solution}

\begin{example}[\markstar]
    证明均值不等式
    \[
        \dfrac{n}{\displaystyle\sum\limits_{i=1}^n\dfrac{1}{x_i}}
        \leq\sqrt[n]{x_1x_2\cdots x_n}
        \leq\dfrac{\displaystyle\sum\limits_{i=1}^n x_i}{n}
        \leq\sqrt{\dfrac{\displaystyle\sum\limits_{i=1}^n x_i^2}{n}}.
    \]
\end{example}
\exampleblank

\begin{solution}
    依次使用加权 AM--GM,AM--GM 与 Cauchy--Schwarz,分别得到调和均值不超过几何均值,几何均值不超过算术均值,算术均值不超过均方根.等号均在 $x_1=\cdots=x_n$ 时成立.
\end{solution}

\begin{example}[\markstar]
    证明 H\"older 不等式: 设 $a_i\geq0$, $x_i\geq0$, $p>1$, $\dfrac{1}{p}+\dfrac{1}{q}=1$, 则
    \[
        \sum\limits_{i=1}^na_ix_i
        \leq\left(\sum\limits_{i=1}^na_i^p\right)^{\frac{1}{p}}
        \left(\sum\limits_{i=1}^nx_i^q\right)^{\frac{1}{q}}.
    \]
\end{example}
\exampleblank

\begin{solution}
    若两范数非零,令 $A=(\sum a_i^p)^{1/p}$,$X=(\sum x_i^q)^{1/q}$.Young 不等式给出
    $a_ix_i/(AX)\leq (a_i/A)^p/p+(x_i/X)^q/q$.求和即得右端为 $1$,从而得到 H\"older 不等式;零范数情形显然.
\end{solution}
